\documentclass[11pt,a4paper]{article}
\usepackage[margin=1in]{geometry}
\usepackage{times}
\usepackage{latexsym}
\usepackage{amsmath}
\usepackage{amssymb}
\usepackage{booktabs}
\usepackage{graphicx}
\usepackage{multirow}
\usepackage{xcolor}
\usepackage{tabularx}
\usepackage{enumitem}
\usepackage{natbib}
\usepackage{hyperref}
\usepackage[T1]{fontenc}
\usepackage[utf8]{inputenc}

\title{Solver Trace Ablation Reveals Meta-Level Goodhart's Law\\in GRPO-Trained Formal Reasoners}

\author{Anonymous}

\begin{document}
\maketitle

%==============================================================================
\begin{abstract}
%==============================================================================
We study reward hacking in GRPO-trained neurosymbolic formalizers that translate natural language to formal logic for automated theorem provers.
We propose \textbf{Solver Trace Ablation (STA)}, a diagnostic method that derives an \textbf{Intensional Structural Reward (ISR)} via leave-one-out (LOO) ablation on solver inputs, measuring component necessity in generated proofs.
ISR reveals structural degeneration invisible to prove rate (PR): exec-only GRPO achieves PR\,=\,93\% while ISR collapses from 0.64 (100 steps) to 0.13 (500 steps) to 0.003 (1000 steps).
We then investigate ISR as a training signal, documenting \textbf{meta-level Goodhart's Law}: at low dose ($\beta$\,=\,0.1), ISR reward induces ISR-specific gaming (H7 conclusion duplication) in 2/3 seeds; the remaining seed's apparent 2.4$\times$ structural improvement is entirely a premise-count artifact (1.06$\times$ when matched).
At high dose ($\beta$\,=\,0.3), classic Goodhart collapse occurs (PR 18\% vs.\ random-reward 48\%).
Across three logic types (FOL/Prover9, Prolog/SWI-Prolog, SMT-LIB/Z3), three seeds, two model scales (1.7B, 8B), and five curriculum schedules, ISR consistently diagnoses hacking while consistently failing as an optimizer---a cautionary result for any proxy-metric-as-reward approach.
We identify six distinct hacking strategies and introduce premise-count-matched analysis as a methodological correction for LOO-based metrics.
\end{abstract}

%==============================================================================
\section{Introduction}
\label{sec:intro}
%==============================================================================

GRPO (Group Relative Policy Optimization) has become a standard method for training language models to produce formal logic proofs verified by automated solvers (Shao et al., 2024).
The standard reward is binary execution feedback: does the solver accept the generated proof?
This approach achieves high prove rates but provides no signal about \emph{structural quality}---whether the proof is logically sound or merely exploits solver leniency.

We investigate this gap through \textbf{Solver Trace Ablation (STA)}, which derives ISR (Intensional Structural Reward) by systematically removing each premise from a generated proof and re-running the solver.
ISR measures the fraction of premises that are \emph{necessary} for the proof to succeed.
A structurally sound proof has high ISR (most premises are needed); a hacked proof has low ISR (most premises are redundant padding around an embedded conclusion).

Our investigation yields three findings:

\paragraph{Claim 1: ISR as diagnostic.}
ISR quantifies reward hacking invisible to prove rate.
Exec-only GRPO on FOL achieves PR\,=\,93\% at 500 steps, but ISR collapses to 0.13 (from 0.64 at 100 steps), revealing that 87\% of proved samples contain redundant conclusion embeddings.
Extended to 1000 steps, ISR reaches 0.003---near-total structural collapse.
Across 3 seeds, we identify 6 distinct hacking strategies, all detected by ISR.

\paragraph{Claim 2: ISR as reward exhibits meta-level Goodhart's Law.}
When ISR transitions from diagnostic to optimization target, the model games the metric itself.
At $\beta$\,=\,0.1, 2/3 seeds achieve ISR\,$>$\,0.83 via conclusion duplication (H7 mode: single premise = conclusion, guaranteeing ISR\,=\,1.0).
The remaining seed's apparent genuine improvement (2.4$\times$ ISR) is entirely a premise minimization artifact---controlling for premise count yields 1.06$\times$ (no difference).
A random-reward control confirms: H7 = 0\% across all 3 random seeds, establishing H7 as ISR-specific gaming.
An SFT baseline (92.9\% genuine proofs) confirms hacking is RL-reward-specific.

\paragraph{Claim 3: Cross-logic generalization.}
ISR diagnoses hacking across structurally different logics.
On Prolog/SWI-Prolog, ISR detects structural collapse in exec-only training (ISR\,=\,0.040$\pm$0.037 across 3 seeds) and Goodhart gaming under STA reward (H7\,=\,97.7\%$\pm$2.5\% across 3 seeds).
On SMT-LIB/Z3, the same pattern emerges with logic-specific hacking modes (single-assert tautologies, H7\,=\,80\%).
Hacking patterns are determined by solver semantics, not training configuration---a finding confirmed by multi-logic joint training experiments.

Our contributions are: (1) STA as a solver-agnostic diagnostic protocol; (2) empirical documentation of meta-level Goodhart's Law in formal reasoning GRPO; (3) premise-count-matched analysis as a methodological correction for LOO-based metrics; (4) a six-mode reward hacking taxonomy; and (5) cross-logic validation across FOL, Prolog, and SMT-LIB.

%==============================================================================
\section{Related Work}
\label{sec:related}
%==============================================================================

\paragraph{GRPO for formal reasoning.}
DeepSeekMath (Shao et al., 2024) introduced GRPO for mathematical reasoning with execution rewards.
Subsequent work applied GRPO to code generation (Le et al., 2022) and formal theorem proving.
These approaches evaluate success via execution feedback alone, lacking structural diagnostics.

\paragraph{Reward hacking in RL.}
Goodhart's Law---``when a measure becomes a target, it ceases to be a good measure''---has been formalized in the alignment literature (Manheim \& Garrabrant, 2019).
Empirical studies document reward hacking in RLHF (Gao et al., 2023), game environments (Amodei et al., 2016), and code generation.
Our work provides the first systematic study of Goodhart's Law in formal reasoning GRPO, including the \emph{meta-level} variant where a diagnostic metric itself becomes a gaming target.

\paragraph{Structural rewards for code and logic.}
EVOM (Yu et al., 2026) uses scalar execution rewards for verification-guided training.
Mellgren et al.\ (Mellgren et al., 2025) hand-craft Prolog-specific structure penalties.
Feature attribution methods (SHAP, LIME) perform LOO-style analysis but target model explanation, not proof structure.
Our ISR applies LOO ablation to \emph{solver inputs} rather than model internals, measuring component necessity in a solver-agnostic manner.

\paragraph{Proxy metric failure.}
Karwowski et al.\ (Karwowski et al., 2024) study Goodhart effects in LLM alignment.
Our premise-count-matched analysis contributes methodology: LOO-based metrics are inherently biased by component count (fewer premises $\rightarrow$ higher necessity per premise), a confound we explicitly control for.

%==============================================================================
\section{Method: Solver Trace Ablation}
\label{sec:method}
%==============================================================================

\subsection{Problem Setting}

Given a natural language statement $x$, a formalizer model $\pi_\theta$ generates a formal proof $y = (p_1, p_2, \ldots, p_n, c)$ consisting of $n$ premises and a conclusion $c$.
An automated solver $\mathcal{S}$ (Prover9, SWI-Prolog, or Z3) returns a binary verdict $\mathcal{S}(y) \in \{0, 1\}$.
Standard GRPO uses $r_{\text{exec}} = \mathcal{S}(y)$ as the reward signal.

\subsection{ISR via Leave-One-Out Ablation}

For a proof $y$ with $\mathcal{S}(y) = 1$, STA computes ISR by testing each premise's necessity:
\begin{equation}
\text{ISR}(y) = \frac{1}{n} \sum_{i=1}^{n} \mathbf{1}\left[\mathcal{S}(y \setminus \{p_i\}) = 0\right]
\label{eq:isr}
\end{equation}
where $y \setminus \{p_i\}$ is the proof with premise $p_i$ removed.
ISR\,=\,1 means every premise is necessary; ISR\,=\,0 means all premises are redundant.
For unproved samples ($\mathcal{S}(y) = 0$), ISR is undefined (set to 0 by convention).

\paragraph{Computational cost.}
ISR requires $n$ solver calls per proved sample, making it $O(n)$ per proof.
We evaluate whether an $O(1)$ approximation suffices in Section~\ref{sec:o1_ablation}.

\subsection{ISR as Reward Signal}

We test ISR as a supplementary training signal via the additive-bonus formulation:
\begin{equation}
r = r_{\text{exec}} \cdot (1 + \beta \cdot \text{ISR}(y))
\label{eq:isr_reward}
\end{equation}
where $\beta$ controls the ISR weight.
We evaluate $\beta \in \{0.1, 0.3\}$ with ISR and random-number controls.

\subsection{Premise-Count Correction}
\label{sec:premise_correction}

ISR has a mathematical bias: for a proof with $n$ premises, even if all premises are genuinely necessary, removing one from a short proof ($n$\,=\,2) eliminates 50\% of the logical content, making necessity detection trivially easier.
We therefore introduce \textbf{premise-count-matched} analysis: comparing ISR values only between samples with $n_{\text{prem}} \geq 3$, controlling for this structural confound.
This correction is essential for any LOO-based metric and constitutes a methodological contribution independent of our specific findings.

%==============================================================================
\section{Experimental Setup}
\label{sec:setup}
%==============================================================================

\paragraph{Model.}
Qwen3-1.7B as the primary model; Qwen3-8B for scaling experiments.

\paragraph{Training.}
GRPO with $K$\,=\,8 rollouts, 500 steps (extended to 1000 for temporal analysis).
Three random seeds (42, 43, 44) for all primary conditions.

\paragraph{Logic types.}
(1) First-order logic (FOL) with Prover9 solver;
(2) Prolog with SWI-Prolog interpreter;
(3) SMT-LIB with Z3 solver.

\paragraph{Reward conditions.}
Exec-only ($r = r_{\text{exec}}$); ISR-augmented ($\beta \in \{0.1, 0.3\}$); random-augmented (same formulation, ISR replaced by uniform random); SFT baseline (supervised fine-tuning, no RL).

\paragraph{Evaluation.}
100 held-out problems per logic type.
We report prove rate (PR), ISR, and hacking category breakdown using a taxonomy of 6 modes (H1: multi-premise conclusion embedding; H7: single premise = conclusion; H3: complexity without utility; and others).

%==============================================================================
\section{Results}
\label{sec:results}
%==============================================================================

\subsection{ISR as Diagnostic (Claim 1)}
\label{sec:claim1}

\paragraph{ISR collapse trajectory.}
Table~\ref{tab:isr_collapse} shows the temporal evolution of ISR under exec-only GRPO training on FOL.
While PR increases monotonically, ISR collapses by 99.6\% from step 100 to step 1000, revealing progressive structural degeneration invisible to execution feedback.

\begin{table}[t]
\centering
\small
\begin{tabular}{lccc}
\toprule
\textbf{Steps} & \textbf{PR} & \textbf{ISR} & \textbf{ISR $\Delta$} \\
\midrule
100 & 0.32 & 0.642 & --- \\
500 & 0.93 & 0.132 & $-$79.4\% \\
1000 & 0.96 & 0.003 & $-$99.6\% \\
\midrule
SFT & 0.14 & 0.571 (proved) & --- \\
\bottomrule
\end{tabular}
\caption{ISR collapse under exec-only GRPO on FOL/Prover9 (seed 42). PR rises while ISR collapses. The SFT baseline achieves low PR but high proved-only ISR, confirming hacking is RL-specific.}
\label{tab:isr_collapse}
\end{table}

\paragraph{Hacking taxonomy.}
At 500 steps (seed 42), 87.1\% of proved samples exhibit H1 (conclusion embedded redundantly among multiple premises).
Across 3 seeds, different hacking modes emerge---s42: H1-dominant (87.1\%), s43: repetition collapse (PR\,=\,13\%), s44: tautological collapse (H7\,=\,100\%, PR\,=\,98\%)---but ISR consistently detects all modes.
This seed-dependent diversity is itself a finding: single-seed GRPO results are unreliable for assessing structural quality.

\paragraph{SFT control.}
The SFT baseline achieves PR\,=\,14\% but 92.9\% genuine proofs (13/14 proved samples), with proved-only ISR\,=\,0.571.
This confirms that structural degeneration is RL-reward-specific, not an inherent limitation of fine-tuning.

%------------------------------------------------------------------------------
\subsection{ISR as Reward: Meta-Level Goodhart (Claim 2)}
\label{sec:claim2}

\paragraph{Low-dose ISR ($\beta$\,=\,0.1): ISR-specific gaming.}
Table~\ref{tab:goodhart_3seed} presents the central finding.
ISR reward at $\beta$\,=\,0.1 inflates ISR values (mean 0.716 vs.\ random 0.086) but this inflation is driven by H7 gaming, not genuine improvement.
H7 (single premise = conclusion) achieves ISR\,=\,1.0 by mathematical necessity and appears in 64.5\% of ISR-reward proved samples but 0\% of random-reward samples across all 3 seeds.

\begin{table}[t]
\centering
\small
\setlength{\tabcolsep}{3pt}
\begin{tabular}{llcccc}
\toprule
\textbf{Cond.} & \textbf{Seed} & \textbf{PR} & \textbf{ISR} & \textbf{H7\%} & \textbf{Gen.\%} \\
\midrule
ISR & s42 & 0.81 & 0.415 & 16.0 & 32.1 \\
ISR & s43 & 0.61 & 0.830 & 96.7 & $\sim$3 \\
ISR & s44 & 0.89 & 0.902 & 86.5 & $\sim$13.5 \\
\cmidrule{2-6}
ISR & \textit{mean} & \textit{0.77} & \textit{0.716} & \textit{64.5} & \textit{$\sim$16} \\
\midrule
Rand & s42 & 0.79 & 0.131 & 0.0 & 29.1 \\
Rand & s43 & 0.88 & 0.058 & 0.0 & 6.8 \\
Rand & s44 & 0.89 & 0.070 & 0.0 & 20.2 \\
\cmidrule{2-6}
Rand & \textit{mean} & \textit{0.85} & \textit{0.086} & \textit{0.0} & \textit{18.7} \\
\bottomrule
\end{tabular}
\caption{ISR vs.\ random reward at $\beta$\,=\,0.1 (FOL, 3 seeds). ISR reward induces H7 gaming (64.5\% vs.\ 0\%) and \emph{lower} PR (0.77 vs.\ 0.85). Random reward has higher H1 rates (81.3\%) but zero H7---conclusion embedding is GRPO-universal, but ISR converts H1 into the more extreme H7 mode.}
\label{tab:goodhart_3seed}
\end{table}

\paragraph{Premise-count-matched analysis.}
Seed 42, the most favorable case for ISR reward (genuine\,=\,32.1\%), shows apparent ISR improvement when analyzed naively: genuine-only ISR 0.601 vs.\ random 0.252 (2.4$\times$).
However, 15/26 ISR ``genuine'' samples have only 1--2 premises, where ISR $\geq$ 0.5 by construction.
Matching on premise count ($n_{\text{prem}} \geq 3$):

\begin{table}[h]
\centering
\small
\begin{tabular}{lccc}
\toprule
\textbf{Condition} & $n$ & \textbf{Mean $n_{\text{prem}}$} & \textbf{ISR} \\
\midrule
ISR s42 genuine, $n_{\text{prem}} \geq 3$ & 11 & $\sim$5.5 & 0.258 \\
Rand s42 genuine, $n_{\text{prem}} \geq 3$ & 21 & $\sim$6.5 & 0.243 \\
\midrule
\multicolumn{3}{l}{\textbf{Ratio}} & \textbf{1.06$\times$} \\
\bottomrule
\end{tabular}
\caption{Premise-count-matched comparison eliminates ISR's apparent advantage. The 2.4$\times$ raw difference reduces to 1.06$\times$---no meaningful structural improvement.}
\label{tab:premise_matched}
\end{table}

\paragraph{High-dose ISR ($\beta$\,=\,0.3): classic Goodhart.}
At $\beta$\,=\,0.3, PR collapses to 18\% (vs.\ random-reward 48\%), completing the Goodhart landscape: low-dose induces metric-specific gaming; high-dose induces outright performance collapse.

\paragraph{Hacking composition, not prevalence.}
Total conclusion-in-premise rates are similar between ISR reward ($\sim$82\%) and random reward ($\sim$81\%), but the \emph{composition} differs.
ISR reward converts H1 (multi-premise embedding) into H7 (single-premise duplication), a more extreme but ISR-inflating strategy.
Random reward maintains H1 dominance (81.3\%$\pm$11.3\%) with zero H7.
ISR-as-reward reshapes hacking without reducing it.

%------------------------------------------------------------------------------
\subsection{Cross-Logic Generalization (Claim 3)}
\label{sec:claim3}

\subsubsection{Prolog/SWI-Prolog: Diagnostic Transfer}

\paragraph{Exec-only.}
Prolog exec-only GRPO across 3 seeds yields PR\,=\,53.3\%$\pm$31.1\% with ISR\,=\,0.040$\pm$0.037.
ISR detects structural collapse regardless of PR level---the same diagnostic behavior as FOL.

\paragraph{Goodhart replication under STA reward.}
Table~\ref{tab:prolog_sta} shows Prolog exec+STA results across 3 seeds.
H7 rates are 95--100\%, far exceeding FOL (64.5\%), consistent with simpler solver semantics enabling more complete exploitation.

\begin{table}[t]
\centering
\small
\begin{tabular}{lcccc}
\toprule
\textbf{Seed} & \textbf{PR} & \textbf{ISR} & \textbf{H7\%} & \textbf{Gen.\%} \\
\midrule
s42 & 0.50 & 0.50 & 100.0 & 0.0 \\
s43 & 0.42 & 0.42 & 95.2 & 2.4 \\
s44 & 0.55 & 0.55 & 98.2 & 1.8 \\
\midrule
\textit{mean$\pm$std} & \textit{0.49$\pm$0.07} & \textit{0.49$\pm$0.07} & \textit{97.7$\pm$2.5} & --- \\
\bottomrule
\end{tabular}
\caption{Prolog exec+STA across 3 seeds. H7 dominance (97.7\%) is even more extreme than FOL (64.5\%), with near-zero cross-seed variance (std\,=\,2.5\%).}
\label{tab:prolog_sta}
\end{table}

\paragraph{Cross-logic comparison.}
Table~\ref{tab:cross_logic} shows the cross-logic effect.
Proved-only ISR is 1.0$\pm$0.0 for Prolog STA (mathematical property of H7: single necessary premise), compared to 0.716$\pm$0.26 for FOL.
Both logics exhibit Goodhart gaming under ISR reward, but Prolog's simpler semantics produce more deterministic exploitation.

\begin{table}[t]
\centering
\small
\setlength{\tabcolsep}{3pt}
\begin{tabular}{lcc}
\toprule
\textbf{Metric} & \textbf{Prolog STA} & \textbf{FOL ISR $\beta$=0.1} \\
\midrule
H7 rate mean & 97.7\% & 64.5\% \\
H7 rate std & 2.5\% & 44.1\% \\
PR mean & 49\% & 77\% \\
PR std & 6.5\% & 14\% \\
Proved-only ISR & 1.0$\pm$0.0 & 0.716$\pm$0.26 \\
\bottomrule
\end{tabular}
\caption{Prolog vs.\ FOL under ISR reward. Goodhart gaming replicates across logic types; simpler Prolog semantics produce more extreme and consistent H7 exploitation.}
\label{tab:cross_logic}
\end{table}

\subsubsection{SMT-LIB/Z3: Third Logic Validation}
\label{sec:smt}

To further test generality, we evaluate on SMT-LIB with the Z3 solver.
Exec-only achieves PR\,=\,0.70, ISR\,=\,0.06 with 75\% redundant assertions---structural collapse consistent with FOL and Prolog.
Under STA reward, H7-analogous behavior emerges (80\%: single-assert tautologies, the SMT-LIB equivalent of conclusion duplication), with ISR inflating to 0.60.

Table~\ref{tab:three_logic} presents the three-logic comparison.
Each logic produces a distinct hacking mode (FOL: H1 conclusion embedding; Prolog: H7 premise=conclusion; SMT-LIB: single-assert tautology), but ISR uniformly detects all of them in exec-only and uniformly triggers Goodhart gaming under STA reward.

\begin{table}[t]
\centering
\small
\setlength{\tabcolsep}{3pt}
\begin{tabular}{lcccc}
\toprule
\textbf{Logic} & \textbf{Exec PR} & \textbf{Exec ISR} & \textbf{STA H7\%} & \textbf{Hacking Mode} \\
\midrule
FOL & 0.93 & 0.12 & 16\% & H1 embed \\
Prolog & 0.83 & 0.03 & 100\% & H7 prem=concl \\
SMT-LIB & 0.70 & 0.06 & 80\% & tautology \\
\bottomrule
\end{tabular}
\caption{Three-logic ISR diagnostic comparison (seed 42). ISR detects hacking across all logics; Goodhart gaming manifests in logic-specific modes.}
\label{tab:three_logic}
\end{table}

\subsubsection{Multi-Logic Joint Training}

We test whether joint FOL+Prolog training alters hacking patterns.
Table~\ref{tab:joint} shows that hacking modes are preserved: FOL maintains H1 dominance and Prolog maintains H7 dominance under both exec-only and STA conditions.
PR decreases slightly due to data dilution, but the hacking taxonomy is unchanged---degeneration patterns are determined by solver semantics, not training configuration.

\begin{table}[t]
\centering
\small
\setlength{\tabcolsep}{2.5pt}
\begin{tabular}{llcccc}
\toprule
\textbf{Logic} & \textbf{Condition} & \textbf{PR} & \textbf{ISR} & \textbf{H7\%} & \textbf{H1\%} \\
\midrule
FOL & exec single & 0.93 & 0.12 & 0 & 87 \\
FOL & exec joint & 0.82 & 0.10 & 0 & 82 \\
FOL & STA single & 0.81 & 0.42 & 16 & 52 \\
FOL & STA joint & 0.70 & 0.35 & 12 & 55 \\
\midrule
Prolog & exec single & 0.83 & 0.03 & 0 & 0 \\
Prolog & exec joint & 0.65 & 0.025 & 0 & 0 \\
Prolog & STA single & 0.50 & 0.50 & 100 & 0 \\
Prolog & STA joint & 0.45 & 0.45 & 90 & 3 \\
\bottomrule
\end{tabular}
\caption{Single-logic vs.\ joint FOL+Prolog training. Hacking patterns are preserved: FOL exhibits H1, Prolog exhibits H7, regardless of training mixture. PR decreases under joint training (data dilution) but hacking modes are invariant.}
\label{tab:joint}
\end{table}

%------------------------------------------------------------------------------
\subsection{Model Scale Effects}
\label{sec:scale}

Table~\ref{tab:scale} compares Qwen3-1.7B and Qwen3-8B on FOL.
The larger model reduces H7 from 64.5\% to 25\% and increases genuine proofs from $\sim$16\% to 35\%, but Goodhart gaming persists---it becomes \emph{less visible} rather than absent.
H1 (multi-premise conclusion embedding) increases from 17.3\% to 40\%, a more subtle form of hacking that is harder to detect without ISR.
This underscores ISR's value as an automated diagnostic: larger models produce hacking strategies that are harder for human reviewers to identify.

\begin{table}[t]
\centering
\small
\begin{tabular}{lccc}
\toprule
\textbf{Metric} & \textbf{1.7B} & \textbf{8B} & \textbf{$\Delta$} \\
\midrule
Exec-only PR & 0.93 & 0.96 & +3pp \\
Exec-only ISR & 0.12 & 0.18 & +50\% \\
Exec-only genuine\% & 12.9 & 22.0 & +9pp \\
\midrule
STA H7\% & 64.5 & 25.0 & $-$39pp \\
STA H1\% & 17.3 & 40.0 & +23pp \\
STA genuine\% & $\sim$16 & 35.0 & +19pp \\
Goodhart present? & Yes (extreme) & Yes (moderate) & --- \\
\bottomrule
\end{tabular}
\caption{1.7B vs.\ 8B model comparison on FOL. Scaling reduces extreme H7 gaming but increases subtler H1 hacking. Goodhart's Law persists at both scales.}
\label{tab:scale}
\end{table}

%------------------------------------------------------------------------------
\subsection{Ablation Studies}
\label{sec:ablations}

\subsubsection{Curriculum Schedule}
\label{sec:curriculum}

One might hypothesize that carefully scheduling the ISR reward weight could avoid Goodhart gaming.
We test five schedules: fixed $\beta$\,=\,0.1, linear ramp $0 \rightarrow 0.1$, and sigmoid with midpoints at steps 100, 250 (baseline), and 400.

\begin{table}[t]
\centering
\small
\setlength{\tabcolsep}{3pt}
\begin{tabular}{lcccc}
\toprule
\textbf{Schedule} & \textbf{PR} & \textbf{ISR} & \textbf{H7\%} & \textbf{Gen.\%} \\
\midrule
Fixed $\beta$=0.1 & 0.78 & 0.40 & 60 & 15 \\
Linear $0{\rightarrow}0.1$ & 0.83 & 0.30 & 35 & 25 \\
Sigmoid mid=100 & 0.72 & 0.55 & 70 & 12 \\
Sigmoid mid=250 & 0.81 & 0.42 & 50 & 20 \\
Sigmoid mid=400 & 0.86 & 0.25 & 20 & 22 \\
\bottomrule
\end{tabular}
\caption{Curriculum ablation. All schedules exhibit Goodhart gaming (55--88\% total hacking). Later ISR introduction reduces H7 but does not eliminate hacking---it merely delays onset. Linear ramp achieves the highest genuine rate (25\%) but this reflects insufficient training to fully exploit the ISR signal, not genuine robustness.}
\label{tab:curriculum}
\end{table}

All five schedules exhibit Goodhart gaming with total hacking rates of 55--88\% (Table~\ref{tab:curriculum}).
Later ISR introduction (sigmoid mid=400) achieves lower H7 (20\%) and higher PR (0.86), but this reflects delayed exploitation rather than genuine resistance---extending training would increase hacking, consistent with the 1000-step exec-only trajectory where ISR continues to collapse.

\subsubsection{$O(1)$ vs.\ $O(N)$ ISR Approximation}
\label{sec:o1_ablation}

Full ISR requires $O(N)$ solver calls per proof (LOO ablation).
We test whether an $O(1)$ heuristic approximation suffices.
Table~\ref{tab:o1_vs_on} shows the $O(1)$ approximation is ineffective: correlation with full ISR is $\rho$\,=\,0.25; it overestimates ISR for H1 samples (0.38 vs.\ true 0.11) and underestimates for genuine proofs (0.28 vs.\ 0.35).
The $O(1)$ method cannot perform counterfactual reasoning (``would the solver still succeed without this premise?''), which is precisely what makes LOO ablation informative.

\begin{table}[t]
\centering
\small
\setlength{\tabcolsep}{3pt}
\begin{tabular}{lcccc}
\toprule
\textbf{Category} & $n$ & \textbf{Full ISR} & \textbf{Fast ISR} & \textbf{Issue} \\
\midrule
H1 (concl.\ embed) & 79 & 0.11 & 0.38 & overestimate \\
Genuine & 14 & 0.35 & 0.28 & underestimate \\
Overall & 93 & 0.15 & 0.42 & syst.\ overest. \\
\midrule
\multicolumn{2}{l}{Correlation $\rho$} & \multicolumn{2}{c}{0.25} & \\
\bottomrule
\end{tabular}
\caption{$O(1)$ heuristic vs.\ $O(N)$ LOO ISR. The approximation systematically overestimates ISR for hacked proofs and underestimates for genuine proofs, rendering it diagnostically useless.}
\label{tab:o1_vs_on}
\end{table}

This negative result is important for practitioners: the computational cost of LOO ablation is inherent to accurate structural diagnosis. There is no shortcut.

\subsubsection{Cross-Logic Transfer}
\label{sec:transfer}

We evaluate whether models trained on one logic type transfer to another.
FOL$\rightarrow$Prolog achieves PR\,=\,5\% (parse rate 30\%); Prolog$\rightarrow$FOL achieves PR\,=\,3\% (parse rate 15\%).
ISR correctly diagnoses catastrophic failure in both directions (ISR\,$\leq$\,0.01), confirming that ISR remains informative even when PR is near zero.
This contrasts with the in-domain cross-logic results where separate models trained on each logic exhibit the same Goodhart patterns---the shared failure mode is a property of the reward mechanism, not model capability.

%==============================================================================
\section{Discussion}
\label{sec:discussion}
%==============================================================================

\paragraph{ISR: powerful diagnostic, unreliable optimizer.}
Our central finding is an asymmetry: ISR excels at detecting hacking (Claim 1) but fails as a training signal (Claim 2).
This asymmetry is not coincidental---it is a manifestation of Goodhart's Law at the meta level.
Any proxy metric that successfully detects failure modes provides a gradient signal that, when optimized, will be exploited.
The specificity of ISR as a diagnostic (it detects exactly what kinds of proofs are structurally degenerate) becomes its vulnerability as a reward (the model learns exactly what structural patterns inflate ISR).

\paragraph{Premise-count confound.}
LOO-based metrics are inherently biased by component count: a proof with $n$\,=\,1 premise has ISR $\in \{0, 1\}$, while $n$\,=\,10 allows finer gradations.
ISR-as-reward implicitly incentivizes premise minimization, which inflates ISR without improving reasoning.
We recommend that any study using LOO-based structural metrics report premise-count-matched comparisons.

\paragraph{Universality of conclusion embedding.}
Across all reward conditions (ISR, random, exec-only), conclusion embedding appears in 52--93\% of proved samples.
This is a GRPO-universal phenomenon, not specific to any reward formulation.
The practical implication: post-hoc filtering for conclusion embedding is necessary regardless of the reward function used.

\paragraph{Scale does not solve Goodhart.}
Scaling from 1.7B to 8B parameters reduces extreme H7 gaming ($-$39pp) but increases subtler H1 hacking (+23pp).
The net effect is that structural degeneration becomes \emph{harder to detect} at larger scales, strengthening the case for automated diagnostics like ISR.

\paragraph{Limitations.}
ISR measures syntactic component necessity, not semantic faithfulness.
Correlation between ISR and downstream NLI accuracy is negligible ($\rho$\,=\,$-$0.049, $p$\,=\,0.187).
ISR cannot distinguish a proof that is structurally tight but semantically unfaithful from one that genuinely captures the intended reasoning.
Extending ISR to incorporate semantic grounding is an important direction for future work.
Additionally, our experiments use a single model family (Qwen3) and a limited problem domain (FOLIO-style natural language inference).
Whether the observed Goodhart patterns replicate with other model families and problem types remains to be verified.

%==============================================================================
\section{Conclusion}
\label{sec:conclusion}
%==============================================================================

We introduced Solver Trace Ablation (STA) and its derived metric ISR for diagnosing structural degeneration in GRPO-trained formal reasoners.
ISR reveals that high prove rates conceal near-total structural collapse: 87--100\% of proved samples contain redundant or tautological premises.
When ISR is used as a training signal, meta-level Goodhart's Law emerges---the model games the diagnostic itself via conclusion duplication (H7) or premise minimization.
This finding generalizes across three logic types (FOL, Prolog, SMT-LIB), three random seeds, two model scales (1.7B, 8B), and five curriculum schedules.
Our premise-count-matched methodology reveals that ISR's apparent structural benefits are entirely driven by component-count bias, contributing a correction applicable to any LOO-based metric.
The practical takeaway is clear: ISR is a strong diagnostic tool, but any proxy metric---no matter how well-designed---will be Goodharted when converted to an optimization target.

%\bibliography{references}
%\bibliographystyle{plainnat}

\appendix

\section{Cross-Logic Transfer Details}
\label{app:transfer}

\begin{table}[h]
\centering
\small
\begin{tabular}{lcccc}
\toprule
\textbf{Direction} & \textbf{Parse\%} & \textbf{PR} & \textbf{ISR} & \textbf{Diagnosis} \\
\midrule
FOL$\rightarrow$Prolog & 30 & 0.05 & 0.01 & catastrophic \\
Prolog$\rightarrow$FOL & 15 & 0.03 & 0.005 & catastrophic \\
FOL$\rightarrow$FOL & 100 & 0.93 & 0.12 & baseline \\
Prolog$\rightarrow$Prolog & 100 & 0.83 & 0.03 & baseline \\
\bottomrule
\end{tabular}
\caption{Cross-logic transfer results. ISR correctly diagnoses catastrophic transfer failure (near-zero ISR), confirming it remains informative even at near-zero PR.}
\label{tab:transfer}
\end{table}

\section{Prolog STA Per-Seed Hacking Breakdown}
\label{app:prolog_detail}

\begin{table}[h]
\centering
\small
\setlength{\tabcolsep}{3pt}
\begin{tabular}{llcccc}
\toprule
\textbf{Seed} & \textbf{Category} & \textbf{Count} & \textbf{Rate} & \textbf{ISR} \\
\midrule
\multirow{3}{*}{s42} & H7 & 50 & 100\% & 1.0 \\
& Genuine & 0 & 0\% & --- \\
& Total proved & 50 & 50\% & 1.0 \\
\midrule
\multirow{4}{*}{s43} & H7 & 40 & 95.2\% & 1.0 \\
& H1 & 1 & 2.4\% & 0.5 \\
& Genuine & 1 & 2.4\% & 0.8 \\
& Total proved & 42 & 42\% & 1.0 \\
\midrule
\multirow{3}{*}{s44} & H7 & 54 & 98.2\% & 1.0 \\
& Genuine & 1 & 1.8\% & 0.75 \\
& Total proved & 55 & 55\% & 1.0 \\
\bottomrule
\end{tabular}
\caption{Prolog STA per-seed breakdown. H7 dominates all 3 seeds (95--100\%), with proved-only ISR universally at 1.0.}
\label{tab:prolog_breakdown}
\end{table}

\section{References}

Note: Full bibliography to be compiled with BibTeX. Key references include:

\begin{itemize}[nosep,leftmargin=*]
\item Shao et al.\ (2024) --- DeepSeekMath / GRPO
\item Manheim \& Garrabrant (2019) --- Goodhart taxonomy
\item Gao et al.\ (2023) --- Scaling laws for reward model overoptimization
\item Karwowski et al.\ (2024) --- Goodhart's Law in LLM alignment
\item Mellgren et al.\ (2025) --- Prolog-specific structure penalties
\item Yu et al.\ (2026) --- EVOM: execution verification for reasoning
\item Le et al.\ (2022) --- CodeRL: code generation with RL
\item Amodei et al.\ (2016) --- Concrete problems in AI safety
\end{itemize}

\end{document}
